This work appeared in the American Journal of Epidemiology \citep{heutonSpatiotemporalForecastingOpioidrelated2024}.  Our Python-based software is available for other researchers to reuse and build upon: \href{https://github.com/tufts-ml/opioid-overdose-models}{\ul{https://github.com/tufts-ml/opioid-overdose-models}}.

\section{Introduction}\label{introduction}

The opioid overdose epidemic in the United States has resulted in over 450,000 deaths during the past eight years, with more than 80,000 fatal opioid-related overdoses during 2022, the highest yet in a single year \citep{fbahmadProvisionalDrugOverdose2023}.
Managing the opioid overdose epidemic requires a constellation of efforts ranging from substance use treatment programs offering medications for opioid use disorder (e.g., methadone, buprenorphine), \citep{mattickBuprenorphineMaintenancePlacebo2014,leshnerMedicationsOpioidUse2019} harm reduction programs with provisions for overdose education and naloxone distribution, and comprehensive mental health and social support services \citep{clarkSystematicReviewCommunity2014,hawkReducingFatalOpioid2015,peiperFentanylTestStrips2019}.
Beyond the provision of harm reduction and healthcare services, it is critical for policymaking to address the ever-evolving substance use environment and plan for targeted interventions.

There has been considerable variation in the availability of different types of opioids and the consequent increase in opioid use disorder and opioid-related fatal overdoses in the past two decades. The current fatal opioid overdose epidemic has been characterized by four waves \citep{ciccaroneTripleWaveEpidemic2019,ciccaroneRiseIllicitFentanyls2021}. In the early 2000s, prescription opioids drove overdoses. Then, heroin-related deaths surged post-2010, followed by a fentanyl spike in 2013 \citep{UnderstandingOpioidOverdose2023}. This culminated with the fourth wave of combined stimulant and fentanyl-related overdoses \citep{ciccaroneRiseIllicitFentanyls2021}. These shifts in supply accompanied changes in social and ecological conditions, impacting substance use behaviors in varying ways across geographic regions \citep{stopkaOpioidrelatedMortalityDynamic2023,friedmanOpioidOverdoseCrisis2020}. Hence, it is critical to examine local spatiotemporal variation in opioid overdose outcomes, identifying the most-impacted areas and predicting future outcomes to inform preemptive public health responses.

A growing body of research \citep{marksMethodologicalApproachesPrediction2021} has explored spatiotemporal variations in the opioid overdose landscape. Yet \emph{forecasting} approaches are in a nascent stage and there are few prediction studies at the population level \citep{borquezFatalOverdosePredicting2022}. Other research focuses on patient-specific risk prediction \citep{tseregounisAssessingOpioidOverdose2021,lo-ciganicIntegratingHumanServices2021,lo-ciganicDevelopingValidatingMachinelearning2022}, assuming access to detailed, person-level demographic and medical history data. However, there are immense challenges in compiling rich datasets for person-level analysis. Most state-level public health authorities may not have access to the data and technical resources to conduct individual predictive modeling. Analyses focused on population-level predictions that solely depend on more readily-available aggregated data have the potential to be more easily adopted by public health authorities with limited resources.

While a number of prior studies have identified historical overdose ``hotspots'' \citep{acharyaExploringCountylevelSpatiotemporal2022,herlandsGaussianProcessSubset2018,luGeographicTemporalTrends2023}, fewer studies have forecasted future spatiotemporal overdose spikes. Research that focuses on hotspots often assumes that identified clusters represent the locations where the highest needs will exist in the future. In our analyses, we show that this assumption does not always hold. Intervention and policies that rely on such findings may be acting on lagged measures of opioid burden, thereby limiting the effectiveness of interventions. Existing research also spans a broad range of spatial and temporal resolutions. In geographic space, studies range from coarser county-level analysis \citep{acharyaExploringCountylevelSpatiotemporal2022}, to finer analyses based on ZIP Codes, census tracts, or census block groups \citep{allenTranslatingPredictiveAnalytics2023}. 
Temporally, studies range in focus across yearly aggregated \citep{acharyaExploringCountylevelSpatiotemporal2022} data, quarterly \citep{heutonZIGP2022}
, or weekly data \citep{ertugrulCASTNetCommunityAttentiveSpatioTemporal2019}.

The overarching goal of our study is to help public health departments make short-term forecasts of future overdose events to enable planning of geographically and temporally targeted interventions that are cognizant of limited resources and needed intervention efficiencies. We focus specifically on development of forecasts at a fine spatial scale (census tracts) at annual intervals, which we selected to match the decision-making needs of public health agencies. In order to understand the role of different forecasting models and evaluation metrics on different communities, our evaluations cover two distinct catchment areas. First, we study Cook County, Illinois, covering over 5 million residents of Chicago and its surroundings, where we forecast death events across 1328 populated census tracts from years 2015-2022 through analysis of publicly available data. Second, we study the state of Massachusetts, where we forecast fatal overdose events across 1620 census tracts representing over 6 million residents from 2001-2021. These locations were selected based on data availability, and to demonstrate the impact of model and metric choice at multiple locations.

To establish best practices for modeling and evaluation, we carefully compare different modeling approaches and performance metrics in each catchment area. We implement a comprehensive set of existing models -- including heuristic baselines, statistical models, \citep{marksMethodologicalApproachesPrediction2021,allenTranslatingPredictiveAnalytics2023,ertugrulCASTNetCommunityAttentiveSpatioTemporal2019,bauerSmallAreaForecasting2023} and neural networks \citep{ertugrulCASTNetCommunityAttentiveSpatioTemporal2019}. We then assess the opioid-related fatal overdose forecasts they produce for both Cook County and Massachusetts at the census-tract-level at annual timeframes. We compute widely-used error-based performance metrics and introduce a new intervention-focused performance metric.
